% ----------------------------------------------------------
% Capítulo 6 — Conclusão
% ----------------------------------------------------------
\chapter{Conclusão}
\label{cap:conclusao}

% ---
\section{Síntese do Trabalho}
% ---

Este Projeto Final de Curso aprimorou a camada de tradução linguística da solução de busca em linguagem natural sobre o acervo cartográfico da Diretoria de Serviço Geográfico (DSG), partindo do protótipo desenvolvido pelo 1º Centro de Geoinformação (1º CGEO). O protótipo demonstrou a viabilidade da abordagem, mas apresentava limitações documentadas pelo próprio cliente: \textit{backend} em Node.js, Saída Estruturada via Instructor e Zod, um único modelo (Phi-4 14B), dicionário de normalização fixo no código, \textit{fallback} por expressões regulares e ausência de um protocolo de avaliação.

A partir da análise crítica do protótipo, o trabalho entregou:

\begin{enumerate}
  \item um \textit{backend} em Python, com FastAPI e LangChain, que substitui o do protótipo e preserva o esquema do banco e o conjunto de parâmetros de busca;
  \item a substituição da Saída Estruturada por \textit{Tool Calling} nativo, em que o modelo decide se chama a ferramenta de busca e com quais argumentos, sem dicionário de normalização e sem exemplos resolvidos no \textit{prompt};
  \item um \textit{dataset} de avaliação com 310 consultas, construído em três camadas (protótipo, autores e geração programática), com gabarito definido por um manual de anotação escrito, leituras múltiplas explícitas para as consultas ambíguas e uma auditoria em três etapas --- checagens automáticas, anotação independente às cegas e adjudicação registrada --- que deixou todas as divergências com decisão justificada;
  \item um \textit{pipeline} de avaliação automatizado e reprodutível, que registra manifesto e traços de cada rodada, recalcula a pontuação contra o gabarito vigente e produz acurácia com intervalo de confiança, precisão, \textit{recall} e F1 por campo, latência e testes de significância entre modelos;
  \item a comparação de quatro modelos abertos executáveis localmente --- Qwen 3 4B, Gemma 4 E4B, Gemma 4 E2B e Mistral Nemo 12B ---, complementada por uma referência paralela com \res{groq}{geral}{nmodelos} modelos de maior porte servidos em nuvem;
  \item um lote de validação independente (lote 1), de \res{lotedados}{geral}{aceitas} consultas com gabarito fixado antes da redação e aceitas por anotação às cegas, sobre o qual as duas abordagens foram comparadas com a especificação da solução e com uma segunda especificação (v2), para separar o efeito da abordagem do efeito da especificação;
  \item a v3 do \textit{Tool Calling}, com ferramentas auxiliares determinísticas e validação com retorno, desenvolvida sobre as 310 consultas e os erros de uma amostra do lote 1 e testada, contra um controle de Saída Estruturada com as mesmas descrições e o mesmo validador, num segundo lote, de \res{lote2dados}{geral}{aceitas} consultas que o desenvolvimento não usou, com a configuração congelada e o plano de análise fixados antes da rodada.
\end{enumerate}

% ---
\section{Retomada dos Objetivos}
% ---

O objetivo geral --- aprimorar a camada de tradução linguística por meio do \textit{Tool Calling} nativo, da avaliação comparativa de modelos locais e de métricas quantitativas --- foi atendido pelos entregáveis definidos na Introdução:

\begin{itemize}
  \item \textbf{Nova arquitetura de \textit{backend}}: descrita no Capítulo~\ref{cap:aprimoramento}, com as camadas de API, orquestração, persistência e avaliação, e com a lógica de consulta ao banco portada do protótipo;
  \item \textbf{\textit{Dataset} de avaliação}: 310 consultas com manual de anotação e registro de auditoria (Capítulo~\ref{cap:metodologia} e Apêndice~\ref{apendice:dataset}); a auditoria mostrou, entre o gabarito e uma anotação independente feita às cegas, concordância quase perfeita nos campos e na detecção de ambiguidade e concordância substancial na decisão de chamar ou não a ferramenta, esta por causa de oito consultas fora do domínio que o gabarito tratava como observacionais, erro corrigido na adjudicação; cada divergência recebeu decisão registrada;
  \item \textbf{\textit{Pipeline} de avaliação}: implementado nos comandos \texttt{pfc-dataset}, \texttt{pfc-auditar}, \texttt{pfc-avaliar} e \texttt{pfc-relatorio}, cobertos por testes automatizados (Capítulo~\ref{cap:aprimoramento});
  \item \textbf{Relatório comparativo e recomendação}: apresentados no Capítulo~\ref{cap:resultados}, com a recomendação de modelo e de abordagem na Seção~\ref{sec:discussao};
  \item \textbf{Lotes de validação e v3}: os dois lotes, com gabarito fixado antes da redação (Seções~\ref{sec:lote-validacao} e~\ref{sec:v3}), e a v3 do \textit{Tool Calling} (Seção~\ref{sec:v3}), implementada no pacote, no avaliador e na API do \textit{backend}, em que é a abordagem padrão (Seção~\ref{sec:harness}); a verificação das conclusões no lote 1 está na Seção~\ref{sec:lote-resultados}, e o teste da v3 no lote 2, na Seção~\ref{sec:v3-resultados}.
\end{itemize}

% ---
\section{Principais Resultados}
% ---

Os resultados do Capítulo~\ref{cap:resultados} permitem responder às perguntas que motivaram o trabalho.

\begin{itemize}
  \item \textbf{Qualidade dos modelos locais}: três dos quatro modelos --- Gemma 4 E4B, Gemma 4 E2B e Qwen 3 4B --- tiveram acurácias estatisticamente indistinguíveis, de \res{principal}{e4b}{acuracia}, \res{principal}{e2b}{acuracia} e \res{principal}{qwen}{acuracia}; o Gemma 4 E4B teve a maior acurácia e o maior F1 (\res{principal}{e4b}{f1}), e o Qwen 3 4B, o menor F1 dos três (\res{principal}{qwen}{f1}), por acrescentar campos não pedidos. O Mistral Nemo 12B, o maior dos quatro, ficou em \res{principal}{mistral}{acuracia}, sobretudo porque deixou de emitir a chamada da ferramenta, que em muitos casos descreveu em texto; o porte do modelo, portanto, não determinou o desempenho.
  \item \textbf{Tool Calling e Saída Estruturada}: com os mesmos modelos e consultas, numa só chamada, o \textit{Tool Calling} não melhorou a extração (a v3, adiante, muda esse quadro). A Saída Estruturada foi mais precisa em três dos quatro modelos locais --- \res{abordagens}{e4b}{acuraciase} contra \res{abordagens}{e4b}{acuraciatc} no Gemma 4 E4B ---, porque o \textit{Tool Calling} dá ao modelo a liberdade de não buscar, e os modelos locais a usaram também em consultas legítimas; o \textit{Tool Calling} só a superou no GPT-OSS 20B, e o efeito dependeu do modelo, não do porte. O ganho consistente do \textit{Tool Calling} foi a recusa: três dos quatro modelos locais e dois dos três modelos em nuvem recusaram todas as \res{principal}{geral}{nFextenso} consultas fora do domínio, e o Gemma 4 E2B, \res{principal}{e2b}{accF} delas, enquanto a Saída Estruturada, na forma avaliada nas 310 e pontuada literalmente, leva toda consulta a uma busca.
  \item \textbf{Lote de validação independente}: num segundo conjunto, o lote 1, de \res{lotedados}{geral}{aceitas} consultas com gabarito fixado antes da redação, o Gemma 4 E4B repetiu a comparação com a especificação da solução (v1) e com a especificação v2, que explicita as convenções do manual e dá às duas abordagens uma forma de recusar. Com a v1, as duas empataram (\res{lote}{tc1}{acuracia} com \textit{Tool Calling} e \res{lote}{se1}{acuracia} com Saída Estruturada; \textit{p}-valor de \res{lote}{tc1se1}{p}), mas o empate soma duas diferenças de sinal contrário: nas consultas do domínio, a Saída Estruturada acertou \res{lote}{se1}{accdominio}, e o \textit{Tool Calling}, \res{lote}{tc1}{accdominio}, sobretudo porque deixou de buscar em \res{lote}{tc1}{falsarecusa} delas, quase sempre para pedir esclarecimento; nas \res{lote}{tc1}{nfora} consultas fora do domínio, o \textit{Tool Calling} recusou todas, e a Saída Estruturada, nenhuma. Pela pontuação do protocolo, em que toda resposta da Saída Estruturada é uma busca, o \textit{Tool Calling} v1 só é mais preciso se mais de \res{lote}{tc1se1}{pstar} das consultas estiverem fora do catálogo; se a aplicação tratasse como não busca a resposta sem nenhum campo do \textit{schema}, a Saída Estruturada v1 recusaria \res{lote}{se1vazio}{recusa} delas, e esse ponto subiria para \res{lote}{tc1se1}{pstarvazio}.
  \item \textbf{Especificação v2}: a v2 levou a falsa recusa do \textit{Tool Calling} a \res{lote}{tc2}{falsarecusa}, mas introduziu erros de forma --- a UF em sigla, o prefixo ``MI'' mantido na \textit{keyword}, o limite acrescentado ---, e a acurácia praticamente não mudou (\res{lote}{tc2}{acuracia}; \textit{p}-valor de \res{lote}{tc1tc2}{p}): o que o \textit{Tool Calling} ganhou ao buscar onde a v1 não buscava, perdeu nas consultas em que as duas versões buscaram, nas quais a acurácia caiu de \res{lote}{tc1tc2}{ambosacc1} para \res{lote}{tc1tc2}{ambosacc2}, e nas subespecificadas, em que passou a buscar com filtros indevidos. Nas 310, a v2 piorou as duas abordagens, embora tenha sido desenhada a partir dos erros nesse conjunto. Com a v2, a ferramenta de recusa do \textit{Tool Calling} teve precisão de \res{lote}{tc2}{recusaprec}, contra \res{lote}{se2}{recusaprec} do campo de recusa da Saída Estruturada, que marcou também \res{lote}{se2}{falsarecusa} das consultas do domínio; a Saída Estruturada, por sua vez, extraiu melhor nas consultas em que as duas abordagens buscaram (\res{lote}{tc2se2}{ambosse} contra \res{lote}{tc2se2}{amboscc}).
  \item \textbf{Arquitetura em duas etapas}: as rodadas do lote 1 permitem simular, sem nova execução, uma combinação das duas vantagens, em que o \textit{Tool Calling} v2 decide se busca ou recusa e, quando busca, valem os parâmetros que a Saída Estruturada v1 extraiu da mesma consulta. A combinação acertou \res{lote}{hib}{acuracia} no lote 1 e \res{base}{hib}{acuracia} nas 310, acima das quatro configurações isoladas nos dois conjuntos. No lote 1, a diferença para a melhor delas, a Saída Estruturada v2, foi de \res{lote}{se2hib}{difacc} pontos percentuais (\textit{p}-valor \res{lote}{se2hib}{p}); nas 310, a combinação ficou acima da Saída Estruturada v1 (\textit{p}-valor de \res{base}{se1hib}{p}), numa diferença que vem inteiramente das consultas fora do domínio. Ela reúne a precisão de recusa do \textit{Tool Calling} v2 (\res{lote}{hib}{recusaprec}, com falsa recusa de \res{lote}{hib}{falsarecusa}) e a extração da Saída Estruturada v1 (\res{lote}{hib}{accdominio} das consultas do domínio, contra \res{lote}{se1}{accdominio} dela sozinha), e supera a Saída Estruturada v1 no lote 1 a partir de \res{lote}{se1hib}{pstar} de consultas fora do domínio, o que torna a escolha praticamente independente dessa proporção. O custo são duas chamadas em série quando há busca: latência mediana de \res{lote}{hib}{latmed}~s na T4, contra \res{lote}{se1}{latmed}~s da Saída Estruturada v1. Foi a melhor combinação das configurações v1 e v2 e, no lote 2, uma das referências secundárias da v3, cuja comparação principal foi com um controle de Saída Estruturada com as mesmas descrições e o mesmo validador.
  \item \textbf{Tool Calling v3, no conjunto de teste}: com ferramentas auxiliares que o modelo chama antes da busca (identificação de UF, município e folha, normalização de código e de escala, resolução de datas pela tabela do manual) e com o retorno de um validador que confere os parâmetros antes de executar, o \textit{Tool Calling} acertou \res{lote2}{tc3}{acuracia} das \res{lote2}{conjunto}{n} consultas do lote 2 (IC 95\%: \res{lote2}{tc3}{ic}). Na comparação principal, fixada antes da rodada, o controle --- a Saída Estruturada com as mesmas descrições e o mesmo validador, mas sem as ferramentas auxiliares --- acertou \res{lote2}{se3}{acuracia} (\textit{p}-valor \res{lote2}{tc3se3}{p}), e a arquitetura em duas etapas simulada com as rodadas do mesmo lote, \res{lote2}{hib}{acuracia}. A vantagem sobre o controle vem toda das consultas do domínio, em que o controle declarou fora do escopo \res{lote2}{se3}{falsarecusa} delas, contra \res{lote2}{tc3}{falsarecusa} do \textit{Tool Calling}, e errou mais parâmetros; o controle acertou mais nas consultas fora do domínio, que recusou todas, e nas subespecificadas (\res{lote2}{se3}{accE} contra \res{lote2}{tc3}{accE}). A partir dos \res{lote2}{tc2}{acuracia} da v2, cada peça da v3 acrescentou acurácia, com diferença significativa, e os degraus medem quanto do resultado é do modelo e quanto é do código: com as descrições corrigidas, ainda numa só chamada, o modelo acertou \res{lote2}{tc3d}{acuracia}; as ferramentas auxiliares, com o usuário simulado, levaram a \res{lote2}{tc3a}{acuracia} (\res{lote2}{tc3dtc3a}{difacc} p.p.), e o retorno do validador, a \res{lote2}{tc3}{acuracia} (\res{lote2}{tc3atc3}{difacc} p.p.). A vantagem do \textit{Tool Calling} já estava na primeira decisão --- a primeira busca ou recusa, depois das ferramentas auxiliares e antes do retorno ---, que acertou \res{lote2}{tc3}{primeira}, contra \res{lote2}{se3}{primeira} no controle (numa comparação exploratória, feita depois do teste, \res{lote2}{tc3se3primeira}{soa} consultas acertadas só pelo \textit{Tool Calling} e \res{lote2}{tc3se3primeira}{sob} só pelo controle; \textit{p}-valor \res{lote2}{tc3se3primeira}{p}), e o retorno do validador a reduziu um pouco: acrescentou mais ao controle (\res{lote2}{se3}{ganhoretorno} p.p.) que ao \textit{Tool Calling} (\res{lote2}{tc3}{ganhoretorno} p.p.) e, das consultas que consertou, deu ao modelo o valor a usar em \res{lote2}{tc3}{consertouvalorpct} no \textit{Tool Calling} e em \res{lote2}{se3}{consertouvalorpct} no controle. A recusa teve precisão de \res{lote2}{tc3}{recusaprec}, mas o seu \textit{recall} caiu para \res{lote2}{tc3}{recusa}, contra \res{lote2}{tc2}{recusa} no \textit{Tool Calling} v2, no degrau de uma chamada, no controle e nas duas etapas; o custo foram cerca de duas chamadas ao modelo por consulta (mediana de \res{lote2}{tc3}{latmed}~s na T4). Nas 310 consultas, que orientaram o desenho e dão um resultado de dentro da amostra, a mesma configuração acertou \res{v3base}{tc3}{acuracia}, e o controle, \res{v3base}{se3}{acuracia}; nos dois conjuntos, o \textit{Tool Calling} v3 ficou acima do controle, e os dois acima das duas etapas e das configurações v1 e v2, mas a ordem entre estas mudou de um conjunto para o outro, e os degraus intermediários não foram executados no desenvolvimento.
  \item \textbf{O protótipo e a solução aprimorada}: o método do protótipo (Phi-4 14B, com dicionário e exemplos) acertou \res{abordagens}{prototipo}{acuracia}, acurácia estatisticamente igual à dos três melhores modelos locais com \textit{Tool Calling} numa só chamada, que são menores, mais de quinze vezes mais rápidos, dispensam o dicionário e recusam o que o catálogo não atende; cada falha passou a ser registrada e classificada, sem \textit{fallback} que a mascarasse. Na solução de uma chamada (v1), o ganho em relação ao protótipo veio, portanto, da troca de modelo e da remoção do dicionário e dos exemplos, e não do \textit{Tool Calling}. A v3 devolve ao sistema código de normalização, mas não como um dicionário aplicado à consulta antes do modelo: como ferramentas que o modelo consulta quando precisa e um validador que confere a saída, com cada intervenção registrada no traço da execução.
  \item \textbf{O que continua difícil}: numa só chamada (v1), datas relativas e ordenação estiveram entre as categorias de menor acurácia para os três modelos locais do patamar superior, e as consultas compostas, que exigem vários campos corretos ao mesmo tempo, não passaram de \res{principal}{e4b}{accC} em nenhum deles. Os modos de falha dominantes foram deixar de chamar a ferramenta --- por descrição da chamada em texto, no Mistral Nemo 12B; por recusa indevida, no Qwen 3 4B; e por pedido de esclarecimento, nos modelos Gemma, em consultas amplas e, no Gemma 4 E4B, também em consultas com datas relativas e com vários critérios --- e, sobretudo no Qwen 3 4B, acrescentar ordenação ou limite não pedidos. Com a v3, essas categorias subiram: numa análise por categoria exploratória, no lote 2, o Gemma 4 E4B acertou \res{lote2}{tc3}{accT} das consultas de tempo relativo, \res{lote2}{tc3}{accO} das de ordenação e \res{lote2}{tc3}{accC} das compostas, contra \res{lote2}{tc1}{accT}, \res{lote2}{tc1}{accO} e \res{lote2}{tc1}{accC} com a v1, com a aritmética de datas passada ao código, que usa a mesma tabela do gabarito; o que a v3 ainda erra está na Seção~\ref{sec:v3-resultados}.
  \item \textbf{\textit{Hardware}}: a qualidade praticamente não mudou entre a estação de referência e a GPU de nuvem; a latência, sim. Na estação de 4~GB, o Gemma 4 E2B teve latência mediana de \res{estacao}{e2b}{latmeds}~s, contra \res{estacao}{qwen}{latmeds}~s do Qwen 3 4B, com acurácia estatisticamente indistinguível (\res{estacao}{e2b}{acuracia} e \res{estacao}{qwen}{acuracia}).
  \item \textbf{Referência em nuvem}: nas mesmas consultas, os modelos servidos pelo Groq acertaram \res{groqcomum}{qwen38}{acuracia} (Qwen 3.8 27B), \res{groqcomum}{gptoss120}{acuracia} (GPT-OSS 120B) e \res{groqcomum}{gptoss20}{acuracia} (GPT-OSS 20B), o que mostra que a tarefa, com a mesma ferramenta, o mesmo \textit{prompt} e o mesmo gabarito, é resolvível por modelos de maior porte. O de maior acurácia, o Qwen 3.8 27B, foi avaliado sem raciocínio, como os modelos locais; a distância indica, portanto, sobretudo a margem de melhoria que modelos de maior capacidade podem trazer.
  \item \textbf{Recomendação}: na solução de uma chamada, recomenda-se, em estações com GPU modesta, o Gemma 4 E2B, pela acurácia estatisticamente indistinguível da dos outros dois melhores modelos locais e pela menor latência nos dois ambientes, com atenção à recusa das consultas fora do domínio; em estações com mais VRAM, o Gemma 4 E4B, que teve a maior acurácia e o maior F1; e o Qwen 3 4B quando predominarem consultas com códigos MI/INOM, tempo relativo ou ordenação. A recomendação de modelo se apoia nas 310 consultas, porque os lotes avaliaram só o Gemma 4 E4B. Quanto à abordagem, recomenda-se o \textit{Tool Calling} v3 com o Gemma 4 E4B, medido só com esse modelo e só na T4: no conjunto de teste, ele superou o controle com Saída Estruturada, a arquitetura em duas etapas e todas as configurações v1 e v2. A troca entre capacidade de recusa e acurácia de extração que as configurações de uma chamada impunham não deixou de existir, mas mudou de forma: o \textit{recall} da recusa caiu para \res{lote2}{tc3}{recusa}, e a v3 custa cerca de duas chamadas ao modelo por consulta. Com o Gemma 4 E2B, na estação de referência, a v3 foi medida só nas 62 consultas P e N, de dentro da amostra (\res{ab}{baseestacaoe2bv3}{acctc} delas, contra \res{ab}{baseestacaoe2bv1}{acctc} da solução v1, com mediana de \res{ab}{baseestacaoe2bv3}{latmedtc}~s); ela é a abordagem padrão da API do \textit{backend}, que permite voltar à solução de uma chamada por configuração (Seção~\ref{sec:harness}). Se for preciso manter uma só chamada, o \textit{Tool Calling} com as descrições corrigidas da v3, sem ferramentas auxiliares nem retorno, acertou \res{lote2}{tc3d}{acuracia} com \res{lote2}{tc3d}{chamadas} chamadas por consulta, em média; se a restrição for a latência, a alternativa é a Saída Estruturada v3, com as mesmas descrições e o mesmo validador, que acertou \res{lote2}{se3}{acuracia}, com latência mediana de \res{lote2}{se3}{latmed}~s na T4, e recusa com menos precisão (falsa recusa de \res{lote2}{se3}{falsarecusa}).
\end{itemize}

% ---
\section{Limitações do Trabalho}
% ---

Os resultados devem ser interpretados à luz das seguintes limitações:

\begin{itemize}
  \item \textbf{Regularidade da camada gerada}: 248 das 310 consultas foram geradas por modelos de frase. A geração garante cobertura sistemática das combinações dos catálogos do domínio, mas essas consultas são linguisticamente mais regulares que as de usuários reais, o que pode favorecer os modelos; por isso os resultados são reportados também por camada (Tabela~\ref{tab:resultados_por_origem}). A camada G responde por \res{principal}{qwen}{nconsultasG} das \res{principal}{geral}{nprincipais} consultas pontuadas, e a acurácia nela superou a das camadas P e N nos três modelos locais do patamar superior --- no Gemma 4 E4B, \res{principal}{e4b}{accorigemG}, contra \res{principal}{e4b}{accorigemN} na camada N e \res{principal}{e4b}{accorigemP} na P --- e nos três modelos em nuvem; as métricas agregadas podem, por isso, superestimar o desempenho em consultas reais. O lote 1, redigido em oito registros de linguagem, confirma a ressalva para o \textit{Tool Calling} v1, cuja falsa recusa passou de \res{base}{tc1}{falsarecusa} nas 310 para \res{lote}{tc1}{falsarecusa} no lote 1 e cuja acurácia nas consultas simples caiu de \res{base}{tc1}{accS} para \res{lote}{tc1}{accS}, mas não para a Saída Estruturada v1, que acertou \res{base}{se1}{accdominio} das consultas do domínio nas 310 e \res{lote}{se1}{accdominio} no lote 1: a superestimação depende da abordagem.
  \item \textbf{Consultas não coletadas de usuários}: nenhuma consulta do \textit{dataset} foi coletada do uso real do catálogo; as camadas P e N foram redigidas pela equipe do protótipo e pelos autores. Os dois lotes também são sintéticos: foram redigidos e anotados por agentes de linguagem de uma mesma família de modelos, e, no lote 1, uma amostra dos erros foi auditada por um único juiz, dessa mesma família; não houve juiz humano. A proporção real de consultas fora do domínio no BDGEx é desconhecida, e dela depende qual das duas abordagens v1 é mais precisa; a vantagem da arquitetura em duas etapas simulada, ao contrário, quase não depende dela, e a da v3 sobre o controle, que vem das consultas do domínio, só se inverteria se a maior parte das consultas estivesse fora dele.
  \item \textbf{Alcance da auditoria}: a anotação independente às cegas foi feita por um modelo de linguagem de propósito geral, e a conferência humana, feita por um dos autores, cobriu uma amostra de 40 consultas; isso verifica a consistência do gabarito com o manual, mas não substitui uma anotação dupla e completa por especialistas do domínio.
  \item \textbf{Ambientes de execução}: as rodadas completas foram executadas em GPU de nuvem, e a estação de referência executou as camadas P e N. A latência depende fortemente do \textit{hardware} e é reportada por ambiente; valores obtidos em outras GPUs, com outra quantidade de VRAM, podem diferir.
  \item \textbf{Linhas de base}: a Saída Estruturada e o método do protótipo foram executados com uma repetição, contra três das rodadas com \textit{Tool Calling} dos modelos locais; o método do protótipo foi reproduzido a partir do código original, com temperatura zero e sem a validação de locais no banco.
  \item \textbf{Alcance dos lotes}: os dois lotes avaliaram um único modelo, o Gemma 4 E4B, numa única GPU, a T4, com uma repetição por configuração, sem medir a variação entre execuções; a recomendação de modelo continua apoiada apenas nas 310 consultas, e a v3 não foi avaliada com outros modelos, entre eles o Gemma 4 E2B, recomendado para as estações com GPU modesta na solução de uma chamada.
  \item \textbf{O código do domínio na v3}: as ferramentas auxiliares e o validador da v3 codificam o manual que define o gabarito --- para datas, usam a mesma tabela que resolve as regras de tempo ---, e o lote 2 sorteou nomes e códigos das mesmas listas que as ferramentas consultam (a lista de municípios do IBGE e o índice do BDGEx). O resultado mede, portanto, o modelo junto com esse código, e o alcance do código em consultas reais, que depende de ele reconhecer nomes e expressões fora dessas listas, não foi medido. Os degraus dizem quanto é de cada um; a decomposição separa só o retorno do validador, porque a primeira decisão, que acertou \res{lote2}{tc3}{primeira}, já vem depois das ferramentas auxiliares, e o que o modelo faz só com a especificação é o degrau de uma chamada (\res{lote2}{tc3d}{acuracia}). Fazem parte do sistema medido também a leitura das chamadas escritas como texto e o pedido ``Responda chamando uma ferramenta'', que o laço envia uma vez quando o modelo responde só com texto. Em \res{lote2}{tc3}{chamadastexto} consultas (\res{lote2}{tc3}{chamadastextopct}), o modelo escreveu ao menos uma chamada no texto da resposta, em vez de usar o canal nativo, e a v3 a interpretou; \res{lote2}{tc3}{chamadastextocertas} delas (\res{lote2}{tc3}{chamadastextocertaspct}) terminaram certas, e, na v1 e na v2, a mesma resposta contaria como não busca. Não se sabe quanto a v3 acertaria sem essa leitura, porque o diálogo seria outro; por isso, ela precisa ir para a produção junto com a v3. O usuário simulado responde sempre que não tem outros detalhes; um usuário real poderia acrescentar critérios, e o efeito do esclarecimento pode ser outro. A v3 foi desenvolvida sobre as 310 consultas e o lote 1, em ciclos sobre os erros de uma amostra deste, e os seus números de desenvolvimento são de dentro da amostra; o lote 2, construído com a mesma receita, mede a generalização para consultas novas, não para outra distribuição de consultas. A análise dos erros e por categoria no lote 2 é exploratória, feita depois do teste, e não mudou a v3. O código precisa ser mantido junto com o catálogo, e a latência da v3 na estação de referência foi medida só nas 62 consultas P e N, com a GPU compartilhada com outros programas.
  \item \textbf{Simulações e atribuições \textit{post hoc}}: a arquitetura em duas etapas e a regra que trata como não busca a resposta sem nenhum campo do \textit{schema} são simulações feitas sobre as rodadas existentes, e não \textit{pipelines} executados; a regra do objeto vazio foi definida depois de ver os dados, e a latência da combinação é a soma das duas chamadas, sem medir \textit{cache} nem sobreposição. As causas atribuídas aos erros de forma da v2 (a lista de siglas sem direção, a instrução de copiar o código exatamente como aparece, a menção do termo proibido e o exemplo de limite) não foram testadas por ablação, e a comparação entre as duas v2 não controla o texto que enquadra a recusa, que difere entre elas. As consequências desses erros na SQL foram simuladas sobre os nomes das UFs e dos municípios, sem o acervo.
  \item \textbf{Critérios de pontuação e de contagem}: na categoria E, a pontuação aceita qualquer não busca, inclusive a recusa, o que favorece o \textit{Tool Calling} v1 e a Saída Estruturada v2, que recusou todas as consultas subespecificadas do lote 1; contadas essas recusas como erro, a Saída Estruturada v2 ficaria em \res{lote}{se2}{accEerro}. No lote 2, o mesmo critério favorece o controle da v3, que não buscou em \res{lote2}{se3}{enao} das subespecificadas, contra \res{lote2}{tc3}{enao} do \textit{Tool Calling} v3, e não explica, portanto, a vantagem da v3. Várias contagens de diagnóstico, como a dos pedidos de esclarecimento e a dos gabaritos discutíveis, dependem do critério adotado, e as análises por registro de linguagem, leitura múltipla e subtipo são \textit{post hoc}, com subgrupos pequenos.
  \item \textbf{Configuração dos modelos}: os modelos locais e o Qwen 3.8 27B foram avaliados sem modo de raciocínio, e os dois GPT-OSS, que não permitem desligá-lo, com o menor esforço de raciocínio que aceitam; todos com temperatura zero e sem exemplos no \textit{prompt}. O desempenho dos modelos locais com raciocínio ativado, e o de todos com \textit{few-shot} ou com \textit{prompts} ajustados a cada modelo, não foi medido.
  \res{principal}{geral}{itemlimitacao}
  \item \textbf{Referência em nuvem}: os modelos do Groq executaram uma repetição, sob os limites de taxa do plano gratuito, e servem apenas como referência de modelos de maior porte; não atendem ao requisito de execução local, e sua latência inclui a rede.
  \item \textbf{Versões dos modelos}: os resultados valem para as versões registradas nos manifestos; o ecossistema de LLMs abertos evolui rapidamente, e o \textit{pipeline} foi construído para que a avaliação possa ser refeita com novos modelos.
\end{itemize}

% ---
\section{Trabalhos Futuros}
% ---

Os resultados e a experiência adquirida apontam as seguintes direções. As duas primeiras, assim como as de interface e implantação, ajuste fino e estudo com usuários, decorrem das exclusões de escopo declaradas no Capítulo~\ref{cap:metodologia}.

\begin{enumerate}
  \item \textbf{Avaliação da recuperação em nível de produto}: medir a fidelidade do conjunto de produtos retornado pela busca (precisão, \textit{recall} e F1 sobre os códigos das folhas), o que exige um gabarito de quais folhas cada consulta deveria devolver.
  \item \textbf{Resolução avançada de nomes geográficos}: resolver os campos \texttt{state} e \texttt{city} em geometrias, para operações espaciais (\texttt{ST\_Intersects}, \texttt{ST\_Within}) contra as \textit{bounding boxes} das folhas, com integração à malha do IBGE e a serviços de geocodificação.
  \item \textbf{\textit{Fallback} para produção}: nas configurações v1 e v2, falhas do modelo são registradas como métrica, sem recuperação por regras determinísticas --- decisão adequada à avaliação; a v3 já recupera parte delas por código, com o retorno do validador, cujas intervenções ficam registradas no traço da execução. Em produção, convém ainda um caminho alternativo, como encaminhar a consulta bruta à busca textual do PostgreSQL.
  \item \textbf{Lacunas do validador da v3}: a leitura exploratória dos erros que restaram no lote 2 mostra o que o validador não confere --- a \textit{keyword} preenchida com região, nome de projeto ou escala --- e o que a tabela do manual não cobre, como algumas expressões de tempo, além das consultas fora do domínio em que o modelo buscou, \res{lote2}{tc3}{contestadasFerrada} delas porque a recusa contestada devolveu ao modelo uma recusa correta. Acrescentar essas regras e rever a contestação é trabalho da mesma natureza, mas o efeito precisa ser medido num conjunto novo, porque o lote 2 já foi usado para identificá-las.
  \item \textbf{A v3 em produção}: a v3 já é a abordagem padrão da API do \textit{backend} (Seção~\ref{sec:harness}), mas sobre a semente sintética; falta executá-la com o banco do acervo --- as rodadas do lote 2 avaliaram os parâmetros, sem executar a SQL ---; medir a sua latência na estação de referência, com, em média, \res{lote2}{tc3}{chamadas} chamadas ao modelo por consulta; levar ao usuário real o pedido de esclarecimento, que aqui e na API tem uma resposta fixa, e avaliá-lo com usuários; e verificar com outros modelos, em especial o Gemma 4 E2B, se a vantagem sobre o controle se mantém.
  \item \textbf{Modos de raciocínio e \textit{prompts}}: avaliar os modelos com raciocínio ativado e com exemplos no \textit{prompt}, medindo o ganho de qualidade contra o custo de latência.
  \item \textbf{Modelos abertos de maior porte em GPU local}: avaliar localmente, numa GPU com 16~GB ou mais de VRAM, modelos de pesos abertos de maior porte, como o GPT-OSS 20B e o Qwen 3.8 27B, que na referência em nuvem acertaram \res{groqcomum}{gptoss20}{acuracia} e \res{groqcomum}{qwen38}{acuracia}, para verificar se a margem observada na nuvem se mantém em execução local, dentro do requisito RNF1.
  \item \textbf{\textit{Dataset} com consultas reais}: ampliar o \textit{dataset} com consultas coletadas de usuários do catálogo, anotadas por especialistas do domínio segundo o manual de anotação, com medida de concordância entre anotadores humanos.
  \item \textbf{Interface e implantação}: desenvolver a interface gráfica, com mapa interativo, e integrar o \textit{pipeline} ao sistema de produção da DSG, com autenticação, controle de acesso e monitoramento.
  \item \textbf{Ajuste fino}: treinar por ajuste fino um modelo de menor porte para o domínio das consultas cartográficas brasileiras, o que pode reduzir o requisito de \textit{hardware}.
  \item \textbf{Estudo com usuários}: avaliar a solução com usuários reais do catálogo, medindo impacto na produtividade e na satisfação.
  \item \textbf{Conversação em múltiplos turnos}: permitir que o usuário refine a busca a partir dos resultados, mantendo o contexto das consultas anteriores.
\end{enumerate}

% ---
\section{Considerações Finais}
% ---

O trabalho mostra que a busca em linguagem natural sobre o acervo cartográfico da DSG pode ser construída com modelos de linguagem abertos, executados localmente e integrados por \textit{Tool Calling}, sem dependência de serviços de nuvem e com o controle sobre os dados que o contexto da DSG e de suas organizações subordinadas exige. Mostra também que o \textit{Tool Calling} não é, por si, garantia de melhor extração: numa só chamada, com os modelos locais avaliados, seu mérito estava na decisão explícita de buscar ou recusar, com precisão de recusa de \res{lote}{tc2}{recusaprec} no lote 1, contra \res{lote}{se2}{recusaprec} do campo de recusa da Saída Estruturada. Usado como mecanismo --- ferramentas auxiliares chamadas no meio da resposta e o retorno de um validador ---, ele superou, num conjunto de teste independente, com o Gemma 4 E4B e em consultas sintéticas, a Saída Estruturada que recebeu as mesmas descrições e o mesmo validador, mas não as ferramentas auxiliares (\res{lote2}{tc3}{acuracia} contra \res{lote2}{se3}{acuracia}); a diferença mede o mecanismo junto com essas ferramentas, que também são código do domínio. A escolha entre as abordagens, ou a forma de combiná-las, deve ser feita com medidas, como as deste trabalho, e não por suposição: o resultado da mesma comparação mudou com a proporção de consultas fora do domínio de cada conjunto, uma especificação desenhada a partir dos erros das 310 piorou nesse mesmo conjunto, e o resultado de teste da v3 vem de um conjunto que o seu desenvolvimento não usou.

Os resultados convergem para uma conclusão que atravessa o trabalho: o diferencial estruturante da solução não é o modelo de linguagem, mas o conhecimento especializado do engenheiro cartógrafo. O modelo contribui com a compreensão da linguagem livre do usuário; a estrutura do domínio vem do especialista. Foi esse conhecimento que definiu o que conta como resposta certa --- o manual de anotação, com as convenções de nomenclatura do INOM e do MI, as escalas do Sistema Cartográfico Nacional, os tipos de produto, os Centros de Geoinformação e as regras de tempo ---, que construiu e auditou os conjuntos de avaliação e que, na v3, virou código: o índice de folhas do acervo, a lista de municípios do IBGE, a normalização de códigos e de escalas, a tabela de datas e o validador que confere cada parâmetro antes da busca. Com o mesmo modelo, esse código levou a acurácia no conjunto de teste de \res{lote2}{tc3d}{acuracia}, só com as descrições corrigidas, a \res{lote2}{tc3}{acuracia}. A experiência da v2 mostra também a forma que esse conhecimento deve ter: escrito como instruções no \textit{prompt} de um modelo pequeno, ele consertou alguns erros e criou outros; escrito como código verificável, testado contra o gabarito e medido num conjunto que o desenvolvimento não viu, ele se somou ao modelo. A ferramenta de IA, sozinha, não substitui a Cartografia; é o engenheiro cartógrafo quem lhe dá a estrutura que a torna útil ao acervo.

Para o cliente, o resultado mais duradouro não é a classificação dos modelos avaliados, que envelhece com o ecossistema, mas o instrumento de avaliação: um \textit{dataset} auditável, com gabarito justificado por regras escritas e divergências adjudicadas, e um \textit{pipeline} reprodutível que permite comparar qualquer novo modelo, \textit{prompt} ou configuração com o mesmo critério.
